\documentclass[12pt,a4paper]{article}

\usepackage{enumitem}

\begin{document}
\begin{titlepage}
    \centering
    \vspace*{2cm}
    {\huge\bfseries Gaussian Blur Filter\par}
    \vspace{1.5cm}
    {\Large\itshape{ Assembler Programming Languages / Project} \par}
    \vspace{2cm}
    {\Large Xawier Słupik, Adam Dyrda\par}
    {\large Student IDs: 317520, 317549\par}
    \vfill
\end{titlepage}

\section{Project Overview}
The primary goal of this project is to implement a high-performance \texttt{Gaussian Blur} image processing algorithm. To achieve optimal execution speed and a user-friendly interface, the project combines high-level application management with low-level assembly optimizations.
\section{Technical Architecture}
\begin{itemize}
    \item Frontend \& UI: Built using C\# .NET WPF to handle the user interface, file input/output, image loading, and thread management.
    \item Blurring Core: Implemented in x64 MASM (Microsoft Macro Assembler) to leverage SIMD vector instructions for optimized image convolution operations.
\end{itemize}

\section{Implementation}
\begin{enumerate}[label=\alph*)]
   \item We begin the implementation with creating the front end application in C\# WPF. The application is designed to load images, manage user interactions, and handle threading for concurrent processing. 
   \item The next thing is to implement the Gaussian Blur algorithm in C\# to understand the idea of the algorithm. And to have a baseline that we can compare our MASM implementation to.
   \item The final step is to implement the highly optimized assembly code to improve the spped of the algorithm even further. 
\end{enumerate}


\end{document}